% Taller / Informe — plantilla con portada en color, listings y diagrama embebido.
\documentclass[12pt]{article}
\usepackage[{{language}}]{babel}
\usepackage[utf8]{inputenc}
\usepackage[T1]{fontenc}
\usepackage{geometry}
\usepackage{xcolor}
\usepackage{hyperref}
\usepackage{enumitem}
\usepackage{fancyhdr}
\usepackage{titlesec}
\usepackage{tcolorbox}
\usepackage{listings}
\usepackage{graphicx}
\usepackage{float}

\geometry{margin=2.5cm}
\hypersetup{colorlinks=true, linkcolor=blue!60!black, urlcolor=blue!60!black, breaklinks=true}

\pagestyle{fancy}
\fancyhf{}
\rhead{{{course}}}
\lhead{{{deliverable}}}
\cfoot{\thepage}

\titleformat{\section}{\large\bfseries\color{blue!50!black}}{\thesection.}{0.5em}{}
\titleformat{\subsection}{\normalsize\bfseries\color{blue!40!black}}{\thesubsection.}{0.5em}{}

% Estilo para listados de código.
\definecolor{codebg}{RGB}{248,248,248}
\definecolor{codeframe}{RGB}{220,220,220}
\lstdefinestyle{codestyle}{
  language=Python,
  basicstyle=\ttfamily\small,
  backgroundcolor=\color{codebg},
  frame=single,
  rulecolor=\color{codeframe},
  breaklines=true,
  showstringspaces=false,
  commentstyle=\color{green!40!black},
  keywordstyle=\color{blue!60!black}\bfseries,
  stringstyle=\color{orange!60!black},
  columns=fullflexible
}

\begin{document}

% ----------------------------------------------------------------------
% Portada
% ----------------------------------------------------------------------
\begin{titlepage}
\centering
{\scshape\LARGE {{deliverable}} \par}
\vspace{1.2cm}
{\huge\bfseries {{title}} \par}
\vspace{0.5cm}
{\Large {{subtitle}} \par}
\vspace{1.6cm}
\begin{tcolorbox}[colback=blue!3, colframe=blue!35!black, width=0.88\textwidth]
\begin{tabular}{ll}
\textbf{Estudiante:} & {{author}} \\
\textbf{Asignatura:} & {{course}} \\
\textbf{Producto:} & {{deliverable}} \\
\textbf{Fecha:} & {{date}} \\
\textbf{Repositorio:} & {\small\url{{{repo_url}}}} \\
\end{tabular}
\end{tcolorbox}
\vfill
\end{titlepage}

% ----------------------------------------------------------------------
% Contenido
% ----------------------------------------------------------------------
\section{Enunciado del ejercicio propuesto}
Describe aquí el objetivo de la actividad y lo que se pide resolver. Resume el
problema en pocas frases y enlaza, si aplica, el repositorio del proyecto en
\url{{{repo_url}}}.

\section{Situación problema}
Explica el contexto y la necesidad concreta que motiva la solución. ¿Qué
dificultad existe hoy y por qué vale la pena automatizarla o resolverla?

\section{Solución algorítmica}
Resume la estrategia de la solución: entradas, salidas, y los pasos generales
del algoritmo.

\begin{tcolorbox}[colback=gray!4, colframe=gray!45, title=Lógica general del algoritmo]
\begin{enumerate}[leftmargin=1.2cm]
\item Inicializar los datos y el estado.
\item Leer la entrada del usuario.
\item Procesar la información según las reglas del problema.
\item Validar los casos límite.
\item Mostrar los resultados de forma organizada.
\end{enumerate}
\end{tcolorbox}

\section{Uso de estructuras de control}
\subsection{Estructuras condicionales}
\begin{itemize}
\item Se usa \texttt{if} para validar la entrada antes de procesarla.
\item Se usa \texttt{if/else} para decidir qué resultado mostrar.
\end{itemize}

\subsection{Estructuras repetitivas}
\begin{itemize}
\item Se usa \texttt{for} para recorrer la colección de datos.
\item Se usa \texttt{while} cuando la cantidad de iteraciones no se conoce de antemano.
\end{itemize}

\section{Diagrama de flujo}
El siguiente diagrama se renderiza automáticamente con \texttt{texforge build}
(no requiere imágenes externas):

\begin{graphviz}[caption=Diagrama de flujo del algoritmo, width=0.55\linewidth]
digraph G {
  rankdir=TB
  node [shape=box, style=rounded, fontname="Helvetica"]
  inicio  [label="Inicio", shape=ellipse]
  leer    [label="Leer entrada"]
  valida  [label="¿Entrada válida?", shape=diamond]
  proc    [label="Procesar datos"]
  error   [label="Mostrar error"]
  salida  [label="Mostrar resultado"]
  fin     [label="Fin", shape=ellipse]

  inicio -> leer -> valida
  valida -> error  [label="no"]
  valida -> proc   [label="sí"]
  proc -> salida -> fin
  error -> fin
}
\end{graphviz}

\section{Código fuente comentado}
El siguiente listado inserta el archivo real de la aplicación. Los comentarios
resaltan el uso de estructuras condicionales y repetitivas.

\lstinputlisting[style=codestyle]{src/example.py}

\section{Reflexión individual}
Comenta qué aprendiste durante el desarrollo, los principales retos y cómo los
resolviste.

\section{Conclusión}
Cierra con los resultados obtenidos y el valor de la solución construida.

\end{document}
